%% 06-theorem-c
\section{Exact Realisation by Path States}
\label{sec:theorem-c}

The two obstructions leave a gap: strictly monotone ultrametric kernels need many
dimensions and cannot be factorised across digits. The construction in this section
fills the gap exactly, and turns out to characterise the whole family.

\begin{theorem}[Path states]
  \label{thm:path-states}
  Let $T$ be a finite rooted tree on node set $V$ with all leaves at depth $n$. For
  every strictly increasing $f : \{0, \dots, n\} \to (0, 1]$ with $f(n) = 1$, set
  \begin{equation}
    \label{eq:amplitudes}
    a_i^{2} = \sqrt{f(i)} - \sqrt{f(i-1)} \quad (0 \le i \le n), \qquad
    \sqrt{f(-1)} := 0,
  \end{equation}
  and define the \emph{path state}
  \begin{equation}
    \label{eq:path-state}
    \pathstate{x} \;=\; \sum_{i=0}^{n} a_i \ket{\anc_i(x)} \;\in\; \C^{V} .
  \end{equation}
  Then $\pathstate{x}$ is a unit vector and, for all leaves $x, y$,
  \begin{equation}
    \label{eq:exactness}
    \fid{\pathstate{x}}{\pathstate{y}} = f\bigl(\lcadepth(x,y)\bigr) \qquad
    \text{exactly.}
  \end{equation}
  Conversely, every kernel that is ultrametric with a strictly monotone profile arises
  this way, and $\Phi_f$ is the unique realisation up to a global isometry and a phase
  on each point \emph{among realisations whose pairwise overlaps are real and
  non-negative}. Dropping that restriction, uniqueness fails
  (\Cref{rem:uniqueness}).
\end{theorem}

\begin{proof}
  Strict monotonicity makes each $a_i^{2}$ in \cref{eq:amplitudes} positive, so the
  $a_i$ are well defined and real. The sum telescopes,
  $\sum_{i=0}^{n} a_i^{2} = \sqrt{f(n)} = 1$, so $\pathstate{x}$ is a unit vector.

  For the overlap, the key structural fact is that two leaves share precisely their
  ancestors down to the depth at which they meet, and no deeper: $\anc_i(x) = \anc_i(y)$
  if and only if $i \le \lcadepth(x,y)$. The basis vectors $\ket{u}$ for $u \in V$ being
  orthonormal,
  \begin{equation}
    \label{eq:overlap}
    \braket{\pathstate{x}}{\pathstate{y}}
    = \sum_{i=0}^{n} a_i^{2} \,\mathbf{1}\bigl[\anc_i(x) = \anc_i(y)\bigr]
    = \sum_{i \le \lcadepth(x,y)} a_i^{2}
    = \sqrt{f\bigl(\lcadepth(x,y)\bigr)} ,
  \end{equation}
  again by telescoping. Squaring the modulus gives \cref{eq:exactness}.

  Two remarks on \cref{eq:overlap}. The cross terms vanish because $\anc_i(x)$ and
  $\anc_j(y)$ sit at different depths when $i \ne j$, so they are distinct nodes and
  hence orthogonal basis vectors. And $a_i^{2} = \sqrt{f(i)} - \sqrt{f(i-1)}$ is forced
  by \cref{eq:overlap} read at successive levels, so $f$ determines the amplitudes.

  For the converse, let $\Psi$ be a realisation whose Gram matrix $G^{\Psi}$ has real
  non-negative entries. Then $G^{\Psi}_{xy} = \sqrt{f(\lcadepth(x,y))}$ exactly, which
  by \cref{eq:overlap} is the Gram matrix of $\Phi_f$. Two families of vectors with
  equal Gram matrices are related by an isometry of their spans, which is the standard
  argument; composing with per-point phases covers the residual freedom that leaves
  $\abs{\cdot}$ unchanged.
\end{proof}

\begin{remark}[How far uniqueness actually goes]
  \label{rem:uniqueness}
  Two weakenings are needed, and the second is not cosmetic.

  First, the converse cannot say ``unique up to relabelling of $V$''. Rotating every
  $\pathstate{x}$ by a fixed unitary, or multiplying each by its own phase, leaves the
  fidelity kernel unchanged while destroying the reading of the basis as tree nodes
  (\texttt{test\_theorem\_c\_uniqueness\_up\_to\_isometry\_and\_phase}). The
  node-labelled form \cref{eq:path-state} is the canonical representative of its
  isometry class, not the only member.

  Second, and less obviously, the positivity hypothesis cannot be dropped. A fidelity
  kernel fixes only the \emph{moduli} $\abs{G_{xy}}$, and these do not determine $G$ up
  to phases: the Bargmann invariants $G_{xy}G_{yz}G_{zx}$ are invariant under
  $\Psi(x) \mapsto e^{i\vartheta_x} U \Psi(x)$ yet are not functions of the moduli.
  Concretely, take $n = 1$, $p = 3$, $f(0) = t = 1/4$ and
  \begin{equation*}
    G^{\Psi} =
    \begin{pmatrix}
      1 & \sqrt{t} & i\sqrt{t} \\
      \sqrt{t} & 1 & \sqrt{t} \\
      -i\sqrt{t} & \sqrt{t} & 1
    \end{pmatrix},
  \end{equation*}
  which is positive semi-definite with spectrum $\{0.134, 1, 1.866\}$ and realises the
  same kernel as $\Phi_f$. Its Bargmann invariant is $-\tfrac{1}{8}i$ against
  $+\tfrac{1}{8}$ for the path state, so no isometry and no choice of phases carries one
  to the other. The correspondence $f \leftrightarrow \Phi_f$ is therefore a bijection
  onto strictly monotone ultrametric kernels only after restricting to realisations with
  non-negative overlaps
  (\texttt{test\_theorem\_c\_converse\_needs\_non\_negative\_overlaps}).
\end{remark}

\begin{remark}[Positive semi-definiteness for free]
  \label{rem:psd}
  Let $A$ be the leaf-by-node incidence matrix with $A_{xu} = a_{\mathrm{depth}(u)}$ if
  $u$ is an ancestor of $x$ and $0$ otherwise. Then \cref{eq:overlap} says
  $\sqrt{K} = A A^{\!\top}$ entrywise, so $K = (A A^{\!\top})^{\circ 2}$ is a Schur
  square of a positive semi-definite matrix and hence positive semi-definite. This
  re-derives, in one line, the classical fact that monotone ultrametric kernels are
  positive semi-definite.
\end{remark}

\subsection{Resources}
\label{subsec:resources}

The path state lives in $\C^{V}$, so it needs $\lceil \log_2 \abs{V} \rceil$ qubits.
Critically, it has only $n+1$ non-zero amplitudes out of $\abs{V}$, so it is a sparse
state. In a level-indexed layout, where the register splits into a depth counter and a
within-level index, preparation is a cascade of $n$ controlled rotations --- one per
level, each splitting the remaining amplitude between ``stop here'' and ``descend'' ---
followed, at each level, by writing the within-level index of $\anc_i(x)$ conditioned on
the depth register. Counting both parts gives $O(n \log \abs{V})$ elementary gates. We
quote that figure throughout; the $O(n)$ rotations alone are not the whole cost.

\begin{table}
  \centering
  \caption{Resources for the path-state encoding, and the \Cref{thm:dimension-bound}
    bound it must respect. Dataset rows use the depth-8 truncation of
    \Cref{sec:experiments}; the WordNet row is the full hierarchy before truncation.}
  \label{tab:resources}
  \begin{tabular}{lrrrr}
    \toprule
    & nodes $\abs{V}$ & qubits & non-zero amplitudes & preparation gates \\
    \midrule
    general tree & $\abs{V}$ & $\lceil \log_2 \abs{V} \rceil$ & $n+1$
      & $O(n \log \abs{V})$ \\
    $\Zpn$ & $\frac{p^{n+1}-1}{p-1}$ & $\approx n \log_2 p$ & $n+1$ & $O(n^{2}\log p)$ \\
    WordNet (full) & $\EOneSourceNodes$ & $17$ & $\EOneSourceHeight + 1$ & --- \\
    WordNet (truncated) & $\PathStateNodes$ & $\PathStateQubits$ & $\EOneHeight + 1$
      & --- \\
    \bottomrule
  \end{tabular}
\end{table}

\Cref{tab:resources} collects the counts. The headline is that the full WordNet noun
hierarchy --- $\EOneSourceNodes$ nodes --- fits in $17$ qubits, with $\EOneSourceHeight
+ 1$ non-zero amplitudes.

\paragraph{Optimality.}
Comparing with \Cref{cor:no-n-qubit}: on $\Zpn$ the construction uses
$\log_2 \abs{V} = n \log_2 p + \log_2 \frac{p}{p-1} + o(1)$ qubits, while the bound
demands at least $n \log_2 p - \log_2 S(p,s)$. The two differ by an additive constant
independent of $n$, so the path-state encoding is optimal up to $O(1)$ qubits. We
resist the stronger phrasing: it is not exactly optimal, and the gap is a genuine, if
small, one.

\subsection{The $p$-adic specialisation}
\label{subsec:padic-specialisation}

On the regular $p$-ary tree, $\lcadepth(x,y) = \min(\vp(x-y), n)$ and the profile can
be tied directly to the $p$-adic absolute value. With $\abs{x-y}_p = p^{-\lcadepth}$ for
$x \ne y$, the geometric profile is
\begin{equation}
  \label{eq:padic-profile}
  f(\lcadepth) = p^{-(n-\lcadepth)s}
  = \Bigl( \frac{p^{-n}}{\abs{x-y}_p} \Bigr)^{\!s} ,
\end{equation}
so the kernel is a decreasing power of the $p$-adic distance, with $f(n) = 1$ filling in
the value at $x = y$ where $\abs{x-y}_p = 0$. The path state is then a level-weighted superposition over the
root-to-leaf path in the $p$-adic tree, with amplitude concentrated near the leaf for
large $s$ and spread towards the root for small $s$. This is the same object that
appears as the state space of continuous-time quantum walks on $p$-adic trees, where
the vertices are sub-balls of a $p$-adic ball
\citep{ZunigaGalindo2024QuantumWalks,ZunigaGalindo2025TwoAdic}, and it connects the
encoding to the $p$-adic quantum formalism of \citet{AnielloEtAl2023PadicQubit}.

\subsection{Arbitrary rooted trees}
\label{subsec:general-trees}

Nothing in the proof of \Cref{thm:path-states} uses regularity. The only structural
input is \cref{eq:overlap}, which needs just that ancestors nest --- true in any rooted
tree. Branching factors may vary between nodes and between levels, and shallow leaves
are handled by the unary padding of \Cref{sec:problem}, which changes no lowest common
ancestor. The construction therefore applies verbatim to WordNet, whose branching
factor ranges from $1$ to $402$, and $\Zpn$ is the homogeneous specialisation rather
than the setting
(\texttt{test\_theorem\_c\_generalises\_to\_a\_non\_homogeneous\_tree}).

\paragraph{Numerical verification.}
Across $(p,n) \in \{(2,3), (3,2), (2,4), (5,2)\}$ and both geometric and linear
profiles, the residual $\max_{x,y} \abs{K - f(\lcadepth)}$ is at most
$\EThreeWorstResidual$, and the induced distance has zero strong-triangle violations
(\texttt{test\_theorem\_c\_path\_states\_realise\_the\_profile\_exactly},
\texttt{test\_theorem\_c\_induced\_distance\_is\_exactly\_ultrametric}).
